\documentclass[12pt,a4paper]{article}
\usepackage{multirow}
\usepackage[utf8]{inputenc}
\usepackage{geometry}
\geometry{margin=0.75in}
\usepackage{mathptmx} % Times New Roman font
\usepackage{helvet}
\usepackage{courier}
\usepackage{amsmath}
\usepackage{amsfonts}
\usepackage{amssymb}
\usepackage{graphicx}
\usepackage{booktabs}
\usepackage{array}
\usepackage{longtable}
\usepackage{url}
\usepackage{xcolor}
\usepackage{hyperref}
\usepackage{subcaption}
\usepackage{float}
\usepackage{algorithm}
\usepackage{algorithmic}
\usepackage{listings}
\usepackage{tcolorbox}
\tcbuselibrary{most}

% Professional color system
\definecolor{primaryblue}{RGB}{31,81,138}
\definecolor{secondaryblue}{RGB}{70,130,180}
\definecolor{accentgreen}{RGB}{40,167,69}
\definecolor{warningred}{RGB}{220,53,69}
\definecolor{highlightgray}{RGB}{248,249,250}
\definecolor{tableheader}{RGB}{233,236,239}

% Professional highlighting commands
\newcommand{\primarytext}[1]{\textcolor{primaryblue}{\textbf{#1}}}
\newcommand{\secondarytext}[1]{\textcolor{secondaryblue}{\textbf{#1}}}
\newcommand{\success}[1]{\textcolor{accentgreen}{\textbf{#1}}}
\newcommand{\warning}[1]{\textcolor{warningred}{\textbf{#1}}}
\newcommand{\highlight}[1]{\colorbox{highlightgray}{#1}}

\lstset{
    basicstyle=\ttfamily\scriptsize,
    breaklines=true,
    frame=single,
    language=Python,
    backgroundcolor=\color{highlightgray}
}

\title{\textbf{Comprehensive Evaluation of Informative Contextual Multi-Armed Bandit Models for Quantum Network Routing}\\
\large{Empirical Assessment of iCMAB Algorithms and Integration Framework for Stochastic Quantum Environments}}

\author{Graduate Research Report\\
Department of Computer Science\\
Rochester Institute of Technology\\
AI/Quantum Computing Research Track}

\date{December 9, 2025}

\begin{document}

\maketitle

\begin{abstract}
This report presents a systematic empirical evaluation of informative contextual multi-armed bandit (iCMAB) algorithms for quantum network routing under stochastic conditions (Attack Rate: 0.25), focusing on stable informative variants (iCEpsilonGreedy, iCPursuit, iCThompsonSampling, iCEpochGreedy, iCEXP4). Through multi-run protocols (3 runs) with stochastic failures, we benchmark performance using Oracle efficiency and validate robustness across configurations. Results show \primarytext{iCEpsilonGreedy leads decisively (86.8\% efficiency)} with excellent baseline performance (93.2\%), followed by \secondarytext{iCPursuit (68.5\% stochastic)} and \success{iCThompsonSampling (66.3\% stochastic)}, establishing a strong foundation for hybrid iCMAB development and deployment in quantum routing under realistic failure conditions.
\end{abstract}

\newpage
\tableofcontents

\newpage

% Executive Summary
\begin{tcolorbox}[colback=highlightgray,colframe=primaryblue,title=\textbf{Executive Summary}]
\begin{itemize}
\item \primarytext{Top Performer}: iCEpsilonGreedy achieves 86.8\% Oracle efficiency (stochastic) with strong baseline 93.2\%, lowest variance (CV = 2.3\%).
\item \secondarytext{Secondary Tier}: iCPursuit (68.5\% stochastic) and iCThompsonSampling (66.3\% stochastic) provide stable alternatives.
\item \success{Robustness}: Coefficient of variation (CV) ranges from 0.6\% to 6.5\% across variants under stochastic conditions.
\item \primarytext{Gap Analysis}: iCEpsilonGreedy achieves minimal Oracle gap (6.4\%) on stochastic workloads, indicating strong adaptive behavior.
\item \warning{Key Finding}: Lower-tier variants (iCEpochGreedy, iCEXP4) achieve only 37\% efficiency, making them unsuitable for critical routing without enhancement.
\end{itemize}
\end{tcolorbox}

\newpage
\section{Introduction}

\subsection{Research Context and Motivation}
Quantum entanglement routing requires adaptive, context-aware decision-making under dynamic network conditions. Informative CMAB (iCMAB) algorithms extend contextual bandits with predictive context modeling to improve stability and efficiency under stochastic failures, providing a principled basis for robust quantum routing.

\subsection{Research Objectives}
\begin{enumerate}
\item \textbf{Performance Characterization}: Identify informative contextual algorithms with superior efficiency and consistency under realistic stochastic failures.
\item \textbf{Robustness Assessment}: Quantify cross-run stability and variance under standardized stochastic conditions.
\item \textbf{Hybrid Integration}: Establish selection criteria and performance thresholds for iCMAB-based hybrid architectures.
\end{enumerate}

\subsection{Contributions}
\begin{itemize}
\item \textbf{Systematic Benchmarking}: Evaluation of five stable iCMAB algorithms under uniform stochastic protocols (Attack Rate: 0.25).
\item \textbf{Statistical Validation}: Multi-run consistency analysis with variance quantification across environments.
\item \textbf{Selection Framework}: Evidence-based guidance for hybrid iCMAB development grounded in empirical data.
\end{itemize}

\section{Theoretical Foundation}

\subsection{Informative Contextual Bandit Formulation}
At step $t$, with context $x_t \in \mathcal{X}$ and informative parameters $\theta$, the expected reward for action $a$ is:
\[
\mu_a(x_t; \theta) = \mathbb{E}[r_{t,a} \mid x_t, \theta]
\]
The goal is to maximize cumulative reward and minimize regret:
\[
\mathrm{Regret}_T = \sum_{t=1}^{T}\left(\max_{a \in \mathcal{A}} \mu_a(x_t;\theta) - \mu_{a_t}(x_t;\theta)\right)
\]

\subsection{Evaluated iCMAB Algorithms}
\begin{itemize}
\item \textbf{iCEpsilonGreedy}: Context-aware epsilon-greedy with informative priors and exploration-exploitation balance.
\item \textbf{iCPursuit}: Informative pursuit with adaptive exploration and contextual confidence updates.
\item \textbf{iCThompsonSampling}: Bayesian sampling with informative posterior updates and uncertainty quantification.
\item \textbf{iCEpochGreedy}: Epoch-based exploration with contextual priors and periodic commitment.
\item \textbf{iCEXP4}: Exponential weights with informative expert signals for expert-based learning.
\end{itemize}

\section{Experimental Methodology}

\subsection{Evaluation Framework}
\begin{itemize}
\item \textbf{Quantum Simulator}: Four-path quantum routing topology with configurable qubit capacities.
\item \textbf{Stochastic Environment}: Attack rate 0.25 to emulate realistic decoherence and network failures.
\item \textbf{Oracle Baseline}: Efficiency computed relative to theoretical upper bound (optimal routing).
\item \textbf{Multi-Run Protocol}: 3 runs per configuration; consistent random seeds for reproducibility.
\item \textbf{Measurement Horizon}: 8000 frames per run for stable convergence assessment.
\end{itemize}

\subsection{Standard Configuration}
\begin{table}[H]
\centering
\caption{Stochastic Evaluation Configuration (Attack Rate: 0.25)}
\begin{tabular}{@{}lll@{}}
\toprule
\textbf{Parameter} & \textbf{Value} & \textbf{Purpose} \\
\midrule
Network Paths & 4 & Realistic routing complexity \\
Qubit Configuration & (8, 10, 8, 9) & Heterogeneous network capacities \\
Frame Horizon & 4000--8000 & Mid-horizon learning window \\
Stochastic Attack Rate & 0.25 (random failures) & Realistic failure simulation \\
Random Seed & Controlled & Reproducibility across runs \\
Metric & Oracle Efficiency (\%) & Normalized performance vs. optimal \\
\bottomrule
\end{tabular}
\end{table}

\section{Results and Analysis}

\begin{tcolorbox}[colback=tableheader,colframe=primaryblue,title=\textbf{Primary iCMAB Performance Summary}]
\subsection{Stochastic Environment Performance (3-Run Aggregate)}
\begin{table}[H]
\centering
\caption{Oracle Efficiency (\%) Under Stochastic Conditions (Attack Rate: 0.25)}
\begin{tabular}{@{}lccccc@{}}
\toprule
\textbf{Algorithm} & \textbf{Run 1} & \textbf{Run 2} & \textbf{Run 3} & \textbf{Avg} & \textbf{CV (\%)} \\
\midrule
\primarytext{iCEpsilonGreedy} & \primarytext{84.0} & \primarytext{88.0} & \primarytext{88.4} & \primarytext{86.8} & \primarytext{2.3} \\
iCPursuit & 70.6 & 62.3 & 72.6 & 68.5 & 6.5 \\
iCThompsonSampling & 64.0 & 66.1 & 68.7 & 66.3 & 2.9 \\
iCEpochGreedy & 37.8 & 37.7 & 37.3 & 37.6 & 0.6 \\
iCEXP4 & 37.4 & 37.1 & 36.7 & 37.1 & 0.8 \\
\bottomrule
\end{tabular}
\end{table}
\end{tcolorbox}

\subsection{Baseline Performance (Optimal Conditions, No Attacks)}
\begin{table}[H]
\centering
\caption{Oracle Efficiency (\%) Under Baseline Conditions (Attack Rate: 0.0)}
\begin{tabular}{@{}lccccc@{}}
\toprule
\textbf{Algorithm} & \textbf{Run 1} & \textbf{Run 2} & \textbf{Run 3} & \textbf{Avg} & \textbf{Gap to Stochastic} \\
\midrule
iCEpsilonGreedy & 92.7 & 93.4 & 93.4 & 93.2 & 6.4\% \\
iCPursuit & 66.2 & 70.4 & 71.1 & 69.2 & 0.7\% \\
iCThompsonSampling & 67.2 & 70.4 & 72.2 & 69.9 & 3.7\% \\
iCEpochGreedy & 39.6 & 39.2 & 39.0 & 39.3 & 1.7\% \\
iCEXP4 & 38.8 & 38.0 & 39.3 & 38.7 & 1.6\% \\
\bottomrule
\end{tabular}
\end{table}

\subsection{Robustness and Consistency Analysis}
\begin{table}[H]
\centering
\caption{iCMAB Robustness Metrics (Stochastic Workload)}
\begin{tabular}{@{}lccccc@{}}
\toprule
\textbf{Algorithm} & \textbf{Mean} & \textbf{Std Dev} & \textbf{CV (\%)} & \textbf{Min} & \textbf{Max} \\
\midrule
\success{iCEpsilonGreedy} & \success{86.8\%} & \success{1.99} & \success{2.3} & \success{84.0} & \success{88.4} \\
iCPursuit & 68.5\% & 4.46 & 6.5 & 62.3 & 72.6 \\
iCThompsonSampling & 66.3\% & 1.92 & 2.9 & 64.0 & 68.7 \\
iCEpochGreedy & 37.6\% & 0.22 & 0.6 & 37.3 & 37.8 \\
iCEXP4 & 37.1\% & 0.29 & 0.8 & 36.7 & 37.4 \\
\bottomrule
\end{tabular}
\end{table}

\section{Performance Tier Classification}

\subsection{Tier 1: High-Performance Algorithms (80\%+ Stochastic Efficiency)}
\textbf{iCEpsilonGreedy}: 86.8\% Oracle efficiency with excellent cross-run consistency (CV = 2.3\%).

This algorithm demonstrates strong adaptive behavior, achieving only a 6.4\% gap between baseline and stochastic conditions. Its performance under realistic failure scenarios makes it the recommended choice for critical quantum routing applications.

\subsection{Tier 2: Moderate-Performance Algorithms (60--75\% Stochastic Efficiency)}
\textbf{iCPursuit} (68.5\%) and \textbf{iCThompsonSampling} (66.3\%) provide secondary deployment options. While iCPursuit shows higher variance (CV = 6.5\%), iCThompsonSampling maintains tighter consistency (CV = 2.9\%). Both degrade gracefully under stochastic conditions with Oracle gaps under 4\%.

\subsection{Tier 3: Lower-Bound Baseline Variants (35--40\% Stochastic Efficiency)}
\textbf{iCEpochGreedy} (37.6\%) and \textbf{iCEXP4} (37.1\%) serve as empirical baselines. Both demonstrate exceptional stability (CV < 1\%) but provide insufficient performance for primary routing decisions. These variants may be useful for comparative analysis or as fallback mechanisms in degraded-mode operations.

\section{Selection Framework for Hybrid Development}

\subsection{Algorithm Recommendation Criteria}
\begin{table}[H]
\centering
\caption{Algorithm Selection Guidelines for Hybrid Integration}
\begin{tabular}{@{}lll@{}}
\toprule
\textbf{Use Case} & \textbf{Primary Choice} & \textbf{Rationale} \\
\midrule
Primary Routing & iCEpsilonGreedy & Highest efficiency (86.8\%), minimal stochastic degradation \\
Uncertainty Quantification & iCThompsonSampling & Bayesian framework with stable variance (CV = 2.9\%) \\
Adaptive Learning & iCPursuit & Context-aware updates, despite higher variance \\
Baseline/Fallback & iCEpochGreely or iCEXP4 & Stable but limited (37\% efficiency) \\
\bottomrule
\end{tabular}
\end{table}

\subsection{iCMAB Hybrid Integration Criteria}
\begin{itemize}
\item \textbf{Efficiency Threshold}: Minimum 80\% Oracle efficiency for Tier 1 deployment in stochastic environments.
\item \textbf{Stability Requirement}: Coefficient of variation \(\leq 3.0\%\) across multi-run configurations.
\item \textbf{Robustness Degradation}: Gap between baseline and stochastic \(<10\%\) to ensure graceful failure mode handling.
\item \textbf{Context Sensitivity}: Demonstrated improvement under context-rich routing scenarios (recommended for hybrid architectures).
\end{itemize}

\section{Limitations and Future Work}

\subsection{Limitations}
\begin{itemize}
\item Evaluation limited to stochastic failure model (Attack Rate: 0.25); full adversarial spectrum remains future work.
\item Three-run configuration provides baseline statistical confidence; extended runs (5, 8, 10) in progress.
\item Network topology fixed at 4 paths; scalability to larger topologies requires additional validation.
\item Lower-tier algorithms (iCEpochGreedy, iCEXP4) may benefit from hyperparameter tuning not performed in this evaluation.
\end{itemize}

\subsection{Future Directions}
\begin{enumerate}
\item Extend evaluation to full adversarial spectrum (Markov, Adaptive, Online-Adaptive attacks).
\item Optimize iCMAB hybrids anchored on iCEpsilonGreedy and iCThompsonSampling with multi-run extended protocols.
\item Scale to larger quantum network topologies and longer horizons (10K--22K frames) for deployment readiness.
\item Investigate hybrid architectures combining iCEpsilonGreedy (primary) with iCThompsonSampling (uncertainty modeling).
\end{enumerate}

\section{Implementation and Reproducibility}

\subsection{Framework Architecture}
\begin{lstlisting}[caption=Core iCMAB Components]
quantum_algorithms.py              # iCMAB implementations
quantum_environment.py             # Stochastic quantum simulation
quantum_evaluator_visualizer.py    # Analysis and visualization
quantum_framework_updated.py       # Integration and orchestration
\end{lstlisting}

\subsection{Reproducibility Guidelines}
\begin{itemize}
\item Standardized parameters across all runs: fixed qubit allocation (8, 10, 8, 9), attack rate 0.25.
\item Controlled random seeds ensure deterministic behavior within statistical variation.
\item Modular framework enables extension to adversarial scenarios and larger topologies.
\item All efficiency metrics computed relative to oracle (optimal routing under identical conditions).
\end{itemize}

\newpage
\section{Conclusions}

\subsection{Key Findings}
\begin{enumerate}
\item \primarytext{iCEpsilonGreedy emerges as the dominant algorithm} with 86.8\% efficiency under stochastic failure conditions and minimal variance (CV = 2.3\%).
\item \secondarytext{Tier 2 alternatives} (iCPursuit, iCThompsonSampling) provide fallback options with graceful degradation (66--69\% efficiency).
\item \success{Empirical baselines} (iCEpochGreely, iCEXP4) demonstrate exceptional stability but insufficient performance for primary deployment.
\item \textbf{Cross-run analysis} validates algorithm selection criteria and establishes performance thresholds for hybrid integration.
\end{enumerate}

\subsection{Recommendations for Hybrid Development}
iCEpsilonGreedy provides the optimal foundation for quantum routing hybrids under realistic stochastic failure scenarios. Integration with uncertainty quantification (iCThompsonSampling) or context-adaptive learning (iCPursuit) can further enhance robustness. Near-term deployment should prioritize iCEpsilonGreedy-anchored architectures with multi-run validation across extended horizons.

\section{Acknowledgments}
Prepared as part of GA work in AI/Quantum Computing at RIT. Evaluation conducted December 9, 2025, focusing on informative contextual bandits for quantum network routing under stochastic conditions.

\end{document}
